\begin{lemma}[Splicing an exact two-color component onto a tree]
Let \(G\) be a finite loopless multigraph with exact edge type \(E\), let
\(S\) be an ordered tree-sequence rooted at \(e_0\), and let \(\mu\) be a
proper partial \(k\)-edge-coloring.  Let \(\alpha\ne\beta\) both belong to
\(\overline\mu(S)\).  Suppose that \(y\notin V(S)\) and that the
\((\alpha,\beta)\)-component through the vertex-node \(y\) contains a
vertex-node belonging to \(V(S)\).

Then there are a vertex \(u\in V(S)\) and a strict ordered extension
\(S'\) of \(S\) such that \(S'\) is obtained from \(S\) by a
\noderef{TashkinovAugmentationSequence}, its last displayed vertex is
\(y\), and every new exact edge and new vertex-node belongs to the
\noderef{TwoColorComponent} through \(y\).  Every new edge has color
\(\alpha\) or \(\beta\), and every new vertex lies outside \(V(S)\).
Moreover, if the new entries have indices
\(|S|\le r<|S'|\), then
\[
 \operatorname{ends}(S'.\operatorname{edge}(r))=
 \begin{cases}
   \{u,S'.\operatorname{vertex}(r+1)\},&r=|S|,\\
   \{S'.\operatorname{vertex}(r),S'.\operatorname{vertex}(r+1)\},
       &r>|S|.
 \end{cases}                                                \tag{1}
\]
Thus the appended entries are the reversed segment from the first
intersection with \(S\) to \(y\), and (1) retains the identity of every
parallel edge rather than retaining only its endpoint pair.
\end{lemma}
\begin{proof}
Put \(K=\noderef{TwoColorComponent}(G,\mu,\alpha,\beta,y)\).  By
\noderef{TwoColorComponentSpanningPath}, properness gives a simple path
\(P\) in the two-color incidence graph whose support is exactly \(K\).
The node \(y\) lies on \(P\), because every connected component contains
its root.  By hypothesis, at least one vertex-node \(x\) of \(P\) has
\(x\in V(S)\).

Orient \(P\) temporarily and let \(t\) be the index of the node \(y\).
Among the finitely many indices of vertex-nodes of \(P\) which belong to
\(V(S)\), choose one, say \(a\), minimizing the ordinary distance to
\(t\).  If it lies after \(t\), reverse \(P\).  Write the subpath from
that chosen node to \(y\) as
\[
 u,e_1,v_1,e_2,v_2,\ldots,e_m,v_m=y.                         \tag{2}
\]
Here \(u\in V(S)\).  The two-color incidence graph
\noderef{TwoColorIncidenceGraph} is bipartite between vertex-nodes and
edge-nodes, so the nodes in (2) alternate in exactly the displayed way.
Since \(y\notin V(S)\), the two ends are different and \(m\ge1\).
Minimality of \(a\) says
\[
                         v_i\notin V(S)\quad(1\le i\le m).    \tag{3}
\]
Indeed, any such \(v_i\in V(S)\) would occur strictly closer to \(y\)
than \(u\).  The simple-path condition says that all nodes in (2) are
distinct.  In particular the vertices \(v_i\) are pairwise distinct and
the symbols \(e_i\) are pairwise distinct individual elements of \(E\).

For each \(i\), adjacency of the three consecutive incidence nodes
\(v_{i-1},e_i,v_i\), where \(v_0=u\), unfolds by
\noderef{TwoColorIncidenceGraph}, and incidence unfolds by
\noderef{Incident}, to
\[
 v_{i-1},v_i\in\operatorname{ends}(e_i),qquad
              \mu(e_i)\in\{\alpha,\beta\}.                  \tag{4}
\]
The two vertices in (4) are distinct, by simplicity and because the graph
is loopless.  By the two-endpoint and loopless semantics of
\noderef{LooplessMultigraph}, two distinct vertices incident with this edge
determine its exact unordered endpoint pair, so
\[
                    \operatorname{ends}(e_i)=\{v_{i-1},v_i\}.\tag{5}
\]
No \(e_i\) is an old displayed edge of \(S\).  To see this with edge
identities intact, suppose \(e_i=S.\operatorname{edge}(r)\).  The
attachment field of the \noderef{TreeSequence} \(S\) says that both
endpoints of this exact edge identity occur in its displayed vertex map;
by the image definition \noderef{TreeSequenceVertexSet}, both therefore
belong to \(V(S)\).  Equation (5) would then put \(v_i\) in \(V(S)\),
contrary to (3).  Thus the old edges and the exact edges in (2) are
disjoint.

Starting with \(S\), append successively the edge--vertex pairs
\[
                         (e_1,v_1),\ldots,(e_m,v_m).          \tag{6}
\]
The old vertex and edge maps are retained literally.  By
\noderef{TreeSequenceVertexSet}, the old vertex set is exactly the image of
the old displayed vertex map.  Hence (3) makes every new \(v_i\) disjoint
from that image, while simple-path injectivity makes the new vertices
pairwise distinct; together these prove injectivity of the enlarged vertex
map.  The preceding paragraph and simple-path injectivity prove injectivity
of the enlarged edge map.  The first-edge field is unchanged.  For the
first new edge use the old attachment vertex \(u\); for every later edge
use the preceding new vertex.  Equation (5) proves each attachment equation.
Consequently (6) defines a \noderef{TreeSequence} \(S'\), makes \(S\) an
exact \noderef{TreeSequenceInitialSegment} of \(S'\), and gives precisely
the two endpoint equations in (1).  It also gives
\(|S'|=|S|+m>|S|\), and its last vertex is \(v_m=y\).

It remains to verify that every step in (6) is a Tashkinov augmentation.
The missing-color witnesses for both \(\alpha\) and \(\beta\) already lie
in \(V(S)\), by the definition \noderef{TreeMissingColors}.  Every partial
tree in (6) retains the old displayed vertex map of \(S\) literally as an
initial segment; the image definition \noderef{TreeSequenceVertexSet}
therefore puts those same witness vertices in every partial tree's vertex
set, while their missing-color facts are properties of \(G\) and \(\mu\)
at the unchanged vertices.  Thus both colors remain tree-missing at every
step.  By (4), the new edge at each step has one of these two colors.  The
one-entry length increase, exact initial-segment equality, edge-color
equality, and retained missing witness are exactly the four clauses of
\noderef{TashkinovAugmentation}.  Concatenating the \(m\) steps in its
reflexive-transitive closure gives the required
\noderef{TashkinovAugmentationSequence}.  Equations (2)--(6) also give all
the stated component-membership and outside-vertex conclusions.
\end{proof}
